\documentclass[12pt,a4paper]{article}
\usepackage[utf8]{inputenc}
\usepackage{amsmath}
\usepackage{amsfonts}
\usepackage{amssymb}
\usepackage{geometry}
\usepackage{cite}
\usepackage{booktabs}
\usepackage{hyperref}
\geometry{margin=1in}

% Removed the complex title formatting that causes scan errors
\title{The Mechanical Prediction for Galactic Rotation Curves: Vacuum Impedance and the Variable Gravitational Constant}
\author{Robin Skavberg}
\date{February 2026}

\begin{document}

\maketitle

\begin{abstract}
This paper presents a mechanical derivation of the gravitational constant $G$ as a function of vacuum restorative stiffness and matter-energy density. By defining $G = \frac{c^2}{8\pi\rho_m}$, we demonstrate that the observed "missing mass" in galactic rotation curves is a result of the vacuum recovering its ground-state tension in low-density environments. We provide a mathematical proof for field-strain satiation and propose a testable prediction regarding gravitational lensing discrepancies between high-density clusters and low-density cosmic voids.
\end{abstract}

\section{The Mechanical Basis of Gravitation}
In standard relativistic astrophysics, $G$ is treated as a universal constant. However, applying a mechanical framework to the vacuum suggests that $G$ represents the coupling efficiency between matter and the elastic fabric of space. We define $c^2$ as the restorative stiffness of the vacuum medium. Matter-energy density ($\rho_m$) acts as the mechanical load, leading to the derivation:
\begin{equation}
G = \frac{c^2}{8\pi\rho_m}
\end{equation}

\section{Impedance, Strain, and Field-Strain Satiation}
We define the vacuum as an elastic medium subject to impedance and strain.

\subsection{Description of Terms}
\begin{itemize}
    \item \textbf{$\rho_m$:} Matter-Energy Density (The mechanical load).
    \item \textbf{$c^2$:} Restorative Stiffness (The intrinsic modulus of the vacuum).
    \item \textbf{$Z_v$:} Vacuum Impedance (Opposition to force).
    \item \textbf{$\epsilon$:} Field Strain (Deformation where $\epsilon \propto 1/G$).
\end{itemize}



\subsection{Mathematical Correlation}
In wave mechanics, impedance ($Z_v$) is the opposition a medium presents to a force. For the vacuum:
\begin{equation}
Z_v \propto \rho_m
\end{equation}
As $\rho_m$ increases, $Z_v$ rises, leading to \textbf{Strain Satiation}. This occurs when the field is mechanically saturated and the coupling $G$ approaches its minimum. Conversely, in the galactic outskirts where $\rho_m \to 0$, $Z_v$ is minimized, and the vacuum asserts its maximum restorative "grip."

\section{Prediction for Galactic Rotation Curves}
The "Dark Matter" effect is reinterpreted here as the recovery of the gravitational constant in low-impedance environments.

\subsection{Suppression vs. Recovery}
Local measurements of $G$ in high-density solar systems are suppressed by local vacuum impedance. As the ambient $\rho_m$ drops in the galactic outskirts:
\begin{enumerate}
    \item The denominator of Eq. (1) decreases.
    \item The value of $G$ increases toward its ground-state maximum.
    \item The increased gravitational coupling provides the necessary centripetal force to maintain stellar orbits without auxiliary mass.
\end{enumerate}



\section{Testable Prediction: Lensing Discrepancy}
Unlike MOND, this model relies on \textbf{density-dependent coupling}. 

\textbf{Experimental Test:} Consider two mass distributions of equal baryonic mass $M$. Distribution $A$ is in a high-density cluster; Distribution $B$ is in a low-density cosmic void. Because $\rho_m(B) < \rho_m(A)$, the effective $G$ for $B$ will be higher. Consequently:
\begin{equation}
\theta_{Void} > \theta_{Cluster} \quad \text{for equal } M
\end{equation}

\newpage
\appendix
\section{Technical Appendix}

\subsection{Vacuum Refractive Index}
The observed gravitational-electromagnetic delay is reconciled via:
\begin{equation}
n_v = \sqrt{1 + \rho_m} \implies v_{\gamma} = \frac{c}{n_v}
\end{equation}

\subsection{Ground-State Restoration Table}
\begin{table}[h!]
\centering
\begin{tabular}{@{}llll@{}}
\toprule
Environment & Density ($\rho_m$) & Effective $G$ & Rotation Effect \\ \midrule
Solar System & High & $G_{std}$ & Newtonian \\
Galactic Disk & Medium & $G > G_{std}$ & Partial Flatness \\
Galactic Halo & Low & $G \to G_{max}$ & Flat Rotation \\
Cosmic Void & $\to 0$ & $G_{max}$ & Max Lensing \\ \bottomrule
\end{tabular}
\end{table}

\begin{thebibliography}{9}
\bibitem{1} Einstein, A. (1916). Annalen der Physik, 35